%%%%%%%%%%%%%%%%%%%%%%%%%%%%%%%%%%%%%%%%%%%%%%%%%%%%%%%%%%%%%%%%%%%%%%%%%%%%%%

%\subsection{An exponential bound on mixing sequences}
%\label{sec: an exponential bound}
%%%%%%%%%%%%%%%%%%%%%%%%%%%%%%%%%%%%%%%%%%%%%%%%%%%%%%%%%%%%%%%%%%%%%%%%%%%%%%%

%\subsection{An exponential bound on mixing sequences}
%\label{sec: an exponential bound}
%%%%%%%%%%%%%%%%%%%%%%%%%%%%%%%%%%%%%%%%%%%%%%%%%%%%%%%%%%%%%%%%%%%%%%%%%%%%%%%

%\subsection{An exponential bound on mixing sequences}
%\label{sec: an exponential bound}
%%%%%%%%%%%%%%%%%%%%%%%%%%%%%%%%%%%%%%%%%%%%%%%%%%%%%%%%%%%%%%%%%%%%%%%%%%%%%%%

%\subsection{An exponential bound on mixing sequences}
%\label{sec: an exponential bound}
%\input{6-3_exponential_bound.tex}

In this sub-section we give a simple upper bound on the number of mixing operations
to perfectly mix $E$ with precision $0$ when $\nofdroplets{E}=n\geq 5$.
This bound is exponential in $n$ --- so it's too weak for our purpose for general $n$ --- but we will
use it only to get a polynomial bound when $n$ is at most $22$
(see Theorem~\ref{thm: perfect mixability polynomial}). 

Let $E$ satisfy Invariant~$(\lambda)$, where $\lambda = I'$ for $n=5,6$ and $\lambda = I$ for $n\geq 7$.
Also, let $A\subseteq E$.
We say that a pair of distinct concentrations in $A$ is \emph{$(\lambda)$-safe} 
if it is $(\lambda)$-safe with respect to $E$, that is, if mixing this pair
preserves Invariant~$(\lambda)$ for $E$.
If there is no $(\lambda)$-safe pair in $A$ then we say that $A$ is \emph{$(\lambda)$-mixed}.
Recall that, since $E$ satisfies Invariant~$(\lambda)$, 
according to the properties established in Sections~\ref{subsec: proof for n greater-equal than 7} and~\ref{subsec: proof for n equal 5},
it is guaranteed that in $E$ there exists a pair of concentrations that is either $(\lambda)$-safe
or near-final;
however, this pair may not be in the set $A$ under consideration. 
(In other words, $A$ being $(\lambda)$-mixed does not mean that all droplets in $A$ have the same concentration
--- it only means that we cannot mix any pairs from $A$ in our perfect-mixing sequence for $E$.)


First, in Lemma~$\ref{lem: E-mixed}$ below we show an upper bound on the number of mixing
operations to $(\lambda)$-mix a subset $A\subseteq E$. Then, we prove Theorem~\ref{thm: perfect mixability upper bound}
that gives an upper bound for the number of mixing operations to perfectly mix $E$ with precision $0$.

%%%%%%%%%%%%%%

\begin{lemma}\label{lem: E-mixed}
	Assume that $E$ satisfies Invariant~$(\lambda)$. 
	If $A \subseteq E$ and $\nofdroplets{A} = k$, 
	then $A$ can be $(\lambda)$-mixed with precision $0$ 
	by a mixing sequence of length at most $(8k^3s(A))^k$.
\end{lemma}

\begin{proof}	
	We prove the lemma by induction with respect to $k$.
	Let $\phi(A)$ denote the number of furthest-apart $(\lambda)$-safe mixing operations needed to $(\lambda)$-mix $A$. We
	use induction to prove that $\phi(A) \le (8k^3s(A))^k$. This will imply the lemma.

	In the base case, when $k\le 1$, we have $\phi(A) = 0$, since $A$ is trivially $(\lambda)$-mixed.
	So for the rest of the proof assume that $k\ge 2$ and that every $A' \subset A$
	with $\nofdroplets{A'} = k'<k$ can be $(\lambda)$-mixed in $\phi(A')\leq(8{k'}^3s(A'))^{k'}$ mixing operations.
	We next show that $A$ can be $(\lambda)$-mixed in at most $(8k^3s(A))^k$ mixing operations.
	
	If $A$ is already $(\lambda)$-mixed then we are done.
	Otherwise, let $x,y\in A$ be the furthest-apart $(\lambda)$-safe pair. 
	If $|x-y|\geq \diameter(A)/k$, we will call pair $x,y$ \emph{acceptable}.
	If this pair $x,y$ is acceptable then mixing $x$ and $y$ decreases 
	$\potential(A)$ at least by a factor of $1-1/2k^3$ (by applying 
	Observation~$\ref{obs: far-appart mix}$ with $\gamma = k$ and $n_A = k$).	
	Then, by Observation~$\ref{obs: far-appart constant mix}$, after at most $2k^3$ 
	such acceptable mixing operations $\potential(A)$ decreases at least by a factor of $\half$, and
	it follows that after at most $2k^3\log{\potential_{A,0}}$ such mixing operations, $A$ becomes $(\lambda)$-mixed.
	However, there may be many steps where there is no acceptable pair.
	To address this, we show below a strategy (basically a divide-and-conquer approach)
	that will bound the number of consecutive steps
	needed for an acceptable pair to appear.
	
	Let $\delta=\diameter(A)$ before any mixing operation.
	Divide the interval $[\min(A), \max(A)]$ into $k$ equal segments 
	such that for at least one segment, say $[l,r]$, no concentration in $A$ lies withing
	the open interval $(l,r)$.
	Split $A$ into $A_1$ and $A_2$ such that $\max(A_1) \leq l$ and $\min(A_2) \geq r$.
	Let $k_i = |A_i|$, for $i=1,2$.
	By our inductive assumption, each $A_i$ can be $(\lambda)$-mixed in
	$\phi(A_i) \le  (8k_i^3s(A_i))^{k_i} < (8k^3s(A))^{k-1}$ mixing operations, respectively.
	Therefore, after at most $\phi(A_1) + \phi(A_2)+1 \le 2(8k^3s(A))^{k-1}$
	mixing operations, one of two things must happen:
	either (i) $A$ is already $(\lambda)$-mixed, or
	(ii) there are $x\in A_1$ and $y\in A_2$ that form an $(\lambda)$-safe pair.
	In option~(ii), $x,y$ are in fact an acceptable pair, by the choice of the segment $[l,r]$.
	As argued in the paragraph above, if we repeat the above strategy, 
	option~(ii) can be repeated no more than $2k^3\log{\potential_{A,0}}$
	times before $A$ becomes $(\lambda)$-mixed. Since $\log{\potential_{A,0}} \le 2s(A)$,
	it follows that
	% 
	\begin{align*}
			\phi(A)
			\;&\leq\; 2k^3\log\potential_{A,0} \cdot 2(8k^3s(A))^{k-1}
			\\
			&\leq\; 4k^3s(A)\cdot 2(8k^3s(A))^{k-1}
			\;\leq\; (8k^3s(A))^k,
	\end{align*}
	%
	completing the inductive step and the proof of the lemma.
\end{proof}

%%%%%%%%%%%%%%

\begin{theorem} \label{thm: perfect mixability upper bound}
	If $|E|=n\geq 5$, then $E$ can be perfectly mixed with a mixing sequence of length at most
	$(8n^3s(E))^n + 144n^3s^2(E)$.
\end{theorem}

\begin{proof}
	Let $E$ satisfy Invariant~$(\lambda)$, where $\lambda = I$ for $n\ge 7$ and $\lambda = I'$ for $n = 5,6$.
	We know that if $E$ satisfies Invariant~$(\lambda)$	then $E$ has a pair of concentrations
	that is either $(\lambda)$-safe or near-final. 
	Therefore, we first $(\lambda)$-mix $E$ using Lemma~$\ref{lem: E-mixed}$, 
	which produces a near-final pair.
	(Such a near-final pair becomes available because, as sown in
	Sections~\ref{subsec: proof for n greater-equal than 7}	and~\ref{subsec: proof for n equal 5},
	$E$ satisfying Invariant~($\lambda$) guarantees the existence of a pair that is either ($\lambda$)-safe
	or near-final, and as $E$ is $(\lambda)$-mixed, the later must hold.)
	We then mix this near-final pair, producing $E$ near-final 
	that can be perfectly mixed using Theorem~$\ref{thm: E near-final}$.	
	The total number of mixing operations in this sequence is at most $(8n^3s(E))^n + 144n^3s^2(E)$.
\end{proof}


In this sub-section we give a simple upper bound on the number of mixing operations
to perfectly mix $E$ with precision $0$ when $\nofdroplets{E}=n\geq 5$.
This bound is exponential in $n$ --- so it's too weak for our purpose for general $n$ --- but we will
use it only to get a polynomial bound when $n$ is at most $22$
(see Theorem~\ref{thm: perfect mixability polynomial}). 

Let $E$ satisfy Invariant~$(\lambda)$, where $\lambda = I'$ for $n=5,6$ and $\lambda = I$ for $n\geq 7$.
Also, let $A\subseteq E$.
We say that a pair of distinct concentrations in $A$ is \emph{$(\lambda)$-safe} 
if it is $(\lambda)$-safe with respect to $E$, that is, if mixing this pair
preserves Invariant~$(\lambda)$ for $E$.
If there is no $(\lambda)$-safe pair in $A$ then we say that $A$ is \emph{$(\lambda)$-mixed}.
Recall that, since $E$ satisfies Invariant~$(\lambda)$, 
according to the properties established in Sections~\ref{subsec: proof for n greater-equal than 7} and~\ref{subsec: proof for n equal 5},
it is guaranteed that in $E$ there exists a pair of concentrations that is either $(\lambda)$-safe
or near-final;
however, this pair may not be in the set $A$ under consideration. 
(In other words, $A$ being $(\lambda)$-mixed does not mean that all droplets in $A$ have the same concentration
--- it only means that we cannot mix any pairs from $A$ in our perfect-mixing sequence for $E$.)


First, in Lemma~$\ref{lem: E-mixed}$ below we show an upper bound on the number of mixing
operations to $(\lambda)$-mix a subset $A\subseteq E$. Then, we prove Theorem~\ref{thm: perfect mixability upper bound}
that gives an upper bound for the number of mixing operations to perfectly mix $E$ with precision $0$.

%%%%%%%%%%%%%%

\begin{lemma}\label{lem: E-mixed}
	Assume that $E$ satisfies Invariant~$(\lambda)$. 
	If $A \subseteq E$ and $\nofdroplets{A} = k$, 
	then $A$ can be $(\lambda)$-mixed with precision $0$ 
	by a mixing sequence of length at most $(8k^3s(A))^k$.
\end{lemma}

\begin{proof}	
	We prove the lemma by induction with respect to $k$.
	Let $\phi(A)$ denote the number of furthest-apart $(\lambda)$-safe mixing operations needed to $(\lambda)$-mix $A$. We
	use induction to prove that $\phi(A) \le (8k^3s(A))^k$. This will imply the lemma.

	In the base case, when $k\le 1$, we have $\phi(A) = 0$, since $A$ is trivially $(\lambda)$-mixed.
	So for the rest of the proof assume that $k\ge 2$ and that every $A' \subset A$
	with $\nofdroplets{A'} = k'<k$ can be $(\lambda)$-mixed in $\phi(A')\leq(8{k'}^3s(A'))^{k'}$ mixing operations.
	We next show that $A$ can be $(\lambda)$-mixed in at most $(8k^3s(A))^k$ mixing operations.
	
	If $A$ is already $(\lambda)$-mixed then we are done.
	Otherwise, let $x,y\in A$ be the furthest-apart $(\lambda)$-safe pair. 
	If $|x-y|\geq \diameter(A)/k$, we will call pair $x,y$ \emph{acceptable}.
	If this pair $x,y$ is acceptable then mixing $x$ and $y$ decreases 
	$\potential(A)$ at least by a factor of $1-1/2k^3$ (by applying 
	Observation~$\ref{obs: far-appart mix}$ with $\gamma = k$ and $n_A = k$).	
	Then, by Observation~$\ref{obs: far-appart constant mix}$, after at most $2k^3$ 
	such acceptable mixing operations $\potential(A)$ decreases at least by a factor of $\half$, and
	it follows that after at most $2k^3\log{\potential_{A,0}}$ such mixing operations, $A$ becomes $(\lambda)$-mixed.
	However, there may be many steps where there is no acceptable pair.
	To address this, we show below a strategy (basically a divide-and-conquer approach)
	that will bound the number of consecutive steps
	needed for an acceptable pair to appear.
	
	Let $\delta=\diameter(A)$ before any mixing operation.
	Divide the interval $[\min(A), \max(A)]$ into $k$ equal segments 
	such that for at least one segment, say $[l,r]$, no concentration in $A$ lies withing
	the open interval $(l,r)$.
	Split $A$ into $A_1$ and $A_2$ such that $\max(A_1) \leq l$ and $\min(A_2) \geq r$.
	Let $k_i = |A_i|$, for $i=1,2$.
	By our inductive assumption, each $A_i$ can be $(\lambda)$-mixed in
	$\phi(A_i) \le  (8k_i^3s(A_i))^{k_i} < (8k^3s(A))^{k-1}$ mixing operations, respectively.
	Therefore, after at most $\phi(A_1) + \phi(A_2)+1 \le 2(8k^3s(A))^{k-1}$
	mixing operations, one of two things must happen:
	either (i) $A$ is already $(\lambda)$-mixed, or
	(ii) there are $x\in A_1$ and $y\in A_2$ that form an $(\lambda)$-safe pair.
	In option~(ii), $x,y$ are in fact an acceptable pair, by the choice of the segment $[l,r]$.
	As argued in the paragraph above, if we repeat the above strategy, 
	option~(ii) can be repeated no more than $2k^3\log{\potential_{A,0}}$
	times before $A$ becomes $(\lambda)$-mixed. Since $\log{\potential_{A,0}} \le 2s(A)$,
	it follows that
	% 
	\begin{align*}
			\phi(A)
			\;&\leq\; 2k^3\log\potential_{A,0} \cdot 2(8k^3s(A))^{k-1}
			\\
			&\leq\; 4k^3s(A)\cdot 2(8k^3s(A))^{k-1}
			\;\leq\; (8k^3s(A))^k,
	\end{align*}
	%
	completing the inductive step and the proof of the lemma.
\end{proof}

%%%%%%%%%%%%%%

\begin{theorem} \label{thm: perfect mixability upper bound}
	If $|E|=n\geq 5$, then $E$ can be perfectly mixed with a mixing sequence of length at most
	$(8n^3s(E))^n + 144n^3s^2(E)$.
\end{theorem}

\begin{proof}
	Let $E$ satisfy Invariant~$(\lambda)$, where $\lambda = I$ for $n\ge 7$ and $\lambda = I'$ for $n = 5,6$.
	We know that if $E$ satisfies Invariant~$(\lambda)$	then $E$ has a pair of concentrations
	that is either $(\lambda)$-safe or near-final. 
	Therefore, we first $(\lambda)$-mix $E$ using Lemma~$\ref{lem: E-mixed}$, 
	which produces a near-final pair.
	(Such a near-final pair becomes available because, as sown in
	Sections~\ref{subsec: proof for n greater-equal than 7}	and~\ref{subsec: proof for n equal 5},
	$E$ satisfying Invariant~($\lambda$) guarantees the existence of a pair that is either ($\lambda$)-safe
	or near-final, and as $E$ is $(\lambda)$-mixed, the later must hold.)
	We then mix this near-final pair, producing $E$ near-final 
	that can be perfectly mixed using Theorem~$\ref{thm: E near-final}$.	
	The total number of mixing operations in this sequence is at most $(8n^3s(E))^n + 144n^3s^2(E)$.
\end{proof}


In this sub-section we give a simple upper bound on the number of mixing operations
to perfectly mix $E$ with precision $0$ when $\nofdroplets{E}=n\geq 5$.
This bound is exponential in $n$ --- so it's too weak for our purpose for general $n$ --- but we will
use it only to get a polynomial bound when $n$ is at most $22$
(see Theorem~\ref{thm: perfect mixability polynomial}). 

Let $E$ satisfy Invariant~$(\lambda)$, where $\lambda = I'$ for $n=5,6$ and $\lambda = I$ for $n\geq 7$.
Also, let $A\subseteq E$.
We say that a pair of distinct concentrations in $A$ is \emph{$(\lambda)$-safe} 
if it is $(\lambda)$-safe with respect to $E$, that is, if mixing this pair
preserves Invariant~$(\lambda)$ for $E$.
If there is no $(\lambda)$-safe pair in $A$ then we say that $A$ is \emph{$(\lambda)$-mixed}.
Recall that, since $E$ satisfies Invariant~$(\lambda)$, 
according to the properties established in Sections~\ref{subsec: proof for n greater-equal than 7} and~\ref{subsec: proof for n equal 5},
it is guaranteed that in $E$ there exists a pair of concentrations that is either $(\lambda)$-safe
or near-final;
however, this pair may not be in the set $A$ under consideration. 
(In other words, $A$ being $(\lambda)$-mixed does not mean that all droplets in $A$ have the same concentration
--- it only means that we cannot mix any pairs from $A$ in our perfect-mixing sequence for $E$.)


First, in Lemma~$\ref{lem: E-mixed}$ below we show an upper bound on the number of mixing
operations to $(\lambda)$-mix a subset $A\subseteq E$. Then, we prove Theorem~\ref{thm: perfect mixability upper bound}
that gives an upper bound for the number of mixing operations to perfectly mix $E$ with precision $0$.

%%%%%%%%%%%%%%

\begin{lemma}\label{lem: E-mixed}
	Assume that $E$ satisfies Invariant~$(\lambda)$. 
	If $A \subseteq E$ and $\nofdroplets{A} = k$, 
	then $A$ can be $(\lambda)$-mixed with precision $0$ 
	by a mixing sequence of length at most $(8k^3s(A))^k$.
\end{lemma}

\begin{proof}	
	We prove the lemma by induction with respect to $k$.
	Let $\phi(A)$ denote the number of furthest-apart $(\lambda)$-safe mixing operations needed to $(\lambda)$-mix $A$. We
	use induction to prove that $\phi(A) \le (8k^3s(A))^k$. This will imply the lemma.

	In the base case, when $k\le 1$, we have $\phi(A) = 0$, since $A$ is trivially $(\lambda)$-mixed.
	So for the rest of the proof assume that $k\ge 2$ and that every $A' \subset A$
	with $\nofdroplets{A'} = k'<k$ can be $(\lambda)$-mixed in $\phi(A')\leq(8{k'}^3s(A'))^{k'}$ mixing operations.
	We next show that $A$ can be $(\lambda)$-mixed in at most $(8k^3s(A))^k$ mixing operations.
	
	If $A$ is already $(\lambda)$-mixed then we are done.
	Otherwise, let $x,y\in A$ be the furthest-apart $(\lambda)$-safe pair. 
	If $|x-y|\geq \diameter(A)/k$, we will call pair $x,y$ \emph{acceptable}.
	If this pair $x,y$ is acceptable then mixing $x$ and $y$ decreases 
	$\potential(A)$ at least by a factor of $1-1/2k^3$ (by applying 
	Observation~$\ref{obs: far-appart mix}$ with $\gamma = k$ and $n_A = k$).	
	Then, by Observation~$\ref{obs: far-appart constant mix}$, after at most $2k^3$ 
	such acceptable mixing operations $\potential(A)$ decreases at least by a factor of $\half$, and
	it follows that after at most $2k^3\log{\potential_{A,0}}$ such mixing operations, $A$ becomes $(\lambda)$-mixed.
	However, there may be many steps where there is no acceptable pair.
	To address this, we show below a strategy (basically a divide-and-conquer approach)
	that will bound the number of consecutive steps
	needed for an acceptable pair to appear.
	
	Let $\delta=\diameter(A)$ before any mixing operation.
	Divide the interval $[\min(A), \max(A)]$ into $k$ equal segments 
	such that for at least one segment, say $[l,r]$, no concentration in $A$ lies withing
	the open interval $(l,r)$.
	Split $A$ into $A_1$ and $A_2$ such that $\max(A_1) \leq l$ and $\min(A_2) \geq r$.
	Let $k_i = |A_i|$, for $i=1,2$.
	By our inductive assumption, each $A_i$ can be $(\lambda)$-mixed in
	$\phi(A_i) \le  (8k_i^3s(A_i))^{k_i} < (8k^3s(A))^{k-1}$ mixing operations, respectively.
	Therefore, after at most $\phi(A_1) + \phi(A_2)+1 \le 2(8k^3s(A))^{k-1}$
	mixing operations, one of two things must happen:
	either (i) $A$ is already $(\lambda)$-mixed, or
	(ii) there are $x\in A_1$ and $y\in A_2$ that form an $(\lambda)$-safe pair.
	In option~(ii), $x,y$ are in fact an acceptable pair, by the choice of the segment $[l,r]$.
	As argued in the paragraph above, if we repeat the above strategy, 
	option~(ii) can be repeated no more than $2k^3\log{\potential_{A,0}}$
	times before $A$ becomes $(\lambda)$-mixed. Since $\log{\potential_{A,0}} \le 2s(A)$,
	it follows that
	% 
	\begin{align*}
			\phi(A)
			\;&\leq\; 2k^3\log\potential_{A,0} \cdot 2(8k^3s(A))^{k-1}
			\\
			&\leq\; 4k^3s(A)\cdot 2(8k^3s(A))^{k-1}
			\;\leq\; (8k^3s(A))^k,
	\end{align*}
	%
	completing the inductive step and the proof of the lemma.
\end{proof}

%%%%%%%%%%%%%%

\begin{theorem} \label{thm: perfect mixability upper bound}
	If $|E|=n\geq 5$, then $E$ can be perfectly mixed with a mixing sequence of length at most
	$(8n^3s(E))^n + 144n^3s^2(E)$.
\end{theorem}

\begin{proof}
	Let $E$ satisfy Invariant~$(\lambda)$, where $\lambda = I$ for $n\ge 7$ and $\lambda = I'$ for $n = 5,6$.
	We know that if $E$ satisfies Invariant~$(\lambda)$	then $E$ has a pair of concentrations
	that is either $(\lambda)$-safe or near-final. 
	Therefore, we first $(\lambda)$-mix $E$ using Lemma~$\ref{lem: E-mixed}$, 
	which produces a near-final pair.
	(Such a near-final pair becomes available because, as sown in
	Sections~\ref{subsec: proof for n greater-equal than 7}	and~\ref{subsec: proof for n equal 5},
	$E$ satisfying Invariant~($\lambda$) guarantees the existence of a pair that is either ($\lambda$)-safe
	or near-final, and as $E$ is $(\lambda)$-mixed, the later must hold.)
	We then mix this near-final pair, producing $E$ near-final 
	that can be perfectly mixed using Theorem~$\ref{thm: E near-final}$.	
	The total number of mixing operations in this sequence is at most $(8n^3s(E))^n + 144n^3s^2(E)$.
\end{proof}


In this sub-section we give a simple upper bound on the number of mixing operations
to perfectly mix $E$ with precision $0$ when $\nofdroplets{E}=n\geq 5$.
This bound is exponential in $n$ --- so it's too weak for our purpose for general $n$ --- but we will
use it only to get a polynomial bound when $n$ is at most $22$
(see Theorem~\ref{thm: perfect mixability polynomial}). 

Let $E$ satisfy Invariant~$(\lambda)$, where $\lambda = I'$ for $n=5,6$ and $\lambda = I$ for $n\geq 7$.
Also, let $A\subseteq E$.
We say that a pair of distinct concentrations in $A$ is \emph{$(\lambda)$-safe} 
if it is $(\lambda)$-safe with respect to $E$, that is, if mixing this pair
preserves Invariant~$(\lambda)$ for $E$.
If there is no $(\lambda)$-safe pair in $A$ then we say that $A$ is \emph{$(\lambda)$-mixed}.
Recall that, since $E$ satisfies Invariant~$(\lambda)$, 
according to the properties established in Sections~\ref{subsec: proof for n greater-equal than 7} and~\ref{subsec: proof for n equal 5},
it is guaranteed that in $E$ there exists a pair of concentrations that is either $(\lambda)$-safe
or near-final;
however, this pair may not be in the set $A$ under consideration. 
(In other words, $A$ being $(\lambda)$-mixed does not mean that all droplets in $A$ have the same concentration
--- it only means that we cannot mix any pairs from $A$ in our perfect-mixing sequence for $E$.)


First, in Lemma~$\ref{lem: E-mixed}$ below we show an upper bound on the number of mixing
operations to $(\lambda)$-mix a subset $A\subseteq E$. Then, we prove Theorem~\ref{thm: perfect mixability upper bound}
that gives an upper bound for the number of mixing operations to perfectly mix $E$ with precision $0$.

%%%%%%%%%%%%%%

\begin{lemma}\label{lem: E-mixed}
	Assume that $E$ satisfies Invariant~$(\lambda)$. 
	If $A \subseteq E$ and $\nofdroplets{A} = k$, 
	then $A$ can be $(\lambda)$-mixed with precision $0$ 
	by a mixing sequence of length at most $(8k^3s(A))^k$.
\end{lemma}

\begin{proof}	
	We prove the lemma by induction with respect to $k$.
	Let $\phi(A)$ denote the number of furthest-apart $(\lambda)$-safe mixing operations needed to $(\lambda)$-mix $A$. We
	use induction to prove that $\phi(A) \le (8k^3s(A))^k$. This will imply the lemma.

	In the base case, when $k\le 1$, we have $\phi(A) = 0$, since $A$ is trivially $(\lambda)$-mixed.
	So for the rest of the proof assume that $k\ge 2$ and that every $A' \subset A$
	with $\nofdroplets{A'} = k'<k$ can be $(\lambda)$-mixed in $\phi(A')\leq(8{k'}^3s(A'))^{k'}$ mixing operations.
	We next show that $A$ can be $(\lambda)$-mixed in at most $(8k^3s(A))^k$ mixing operations.
	
	If $A$ is already $(\lambda)$-mixed then we are done.
	Otherwise, let $x,y\in A$ be the furthest-apart $(\lambda)$-safe pair. 
	If $|x-y|\geq \diameter(A)/k$, we will call pair $x,y$ \emph{acceptable}.
	If this pair $x,y$ is acceptable then mixing $x$ and $y$ decreases 
	$\potential(A)$ at least by a factor of $1-1/2k^3$ (by applying 
	Observation~$\ref{obs: far-appart mix}$ with $\gamma = k$ and $n_A = k$).	
	Then, by Observation~$\ref{obs: far-appart constant mix}$, after at most $2k^3$ 
	such acceptable mixing operations $\potential(A)$ decreases at least by a factor of $\half$, and
	it follows that after at most $2k^3\log{\potential_{A,0}}$ such mixing operations, $A$ becomes $(\lambda)$-mixed.
	However, there may be many steps where there is no acceptable pair.
	To address this, we show below a strategy (basically a divide-and-conquer approach)
	that will bound the number of consecutive steps
	needed for an acceptable pair to appear.
	
	Let $\delta=\diameter(A)$ before any mixing operation.
	Divide the interval $[\min(A), \max(A)]$ into $k$ equal segments 
	such that for at least one segment, say $[l,r]$, no concentration in $A$ lies withing
	the open interval $(l,r)$.
	Split $A$ into $A_1$ and $A_2$ such that $\max(A_1) \leq l$ and $\min(A_2) \geq r$.
	Let $k_i = |A_i|$, for $i=1,2$.
	By our inductive assumption, each $A_i$ can be $(\lambda)$-mixed in
	$\phi(A_i) \le  (8k_i^3s(A_i))^{k_i} < (8k^3s(A))^{k-1}$ mixing operations, respectively.
	Therefore, after at most $\phi(A_1) + \phi(A_2)+1 \le 2(8k^3s(A))^{k-1}$
	mixing operations, one of two things must happen:
	either (i) $A$ is already $(\lambda)$-mixed, or
	(ii) there are $x\in A_1$ and $y\in A_2$ that form an $(\lambda)$-safe pair.
	In option~(ii), $x,y$ are in fact an acceptable pair, by the choice of the segment $[l,r]$.
	As argued in the paragraph above, if we repeat the above strategy, 
	option~(ii) can be repeated no more than $2k^3\log{\potential_{A,0}}$
	times before $A$ becomes $(\lambda)$-mixed. Since $\log{\potential_{A,0}} \le 2s(A)$,
	it follows that
	% 
	\begin{align*}
			\phi(A)
			\;&\leq\; 2k^3\log\potential_{A,0} \cdot 2(8k^3s(A))^{k-1}
			\\
			&\leq\; 4k^3s(A)\cdot 2(8k^3s(A))^{k-1}
			\;\leq\; (8k^3s(A))^k,
	\end{align*}
	%
	completing the inductive step and the proof of the lemma.
\end{proof}

%%%%%%%%%%%%%%

\begin{theorem} \label{thm: perfect mixability upper bound}
	If $|E|=n\geq 5$, then $E$ can be perfectly mixed with a mixing sequence of length at most
	$(8n^3s(E))^n + 144n^3s^2(E)$.
\end{theorem}

\begin{proof}
	Let $E$ satisfy Invariant~$(\lambda)$, where $\lambda = I$ for $n\ge 7$ and $\lambda = I'$ for $n = 5,6$.
	We know that if $E$ satisfies Invariant~$(\lambda)$	then $E$ has a pair of concentrations
	that is either $(\lambda)$-safe or near-final. 
	Therefore, we first $(\lambda)$-mix $E$ using Lemma~$\ref{lem: E-mixed}$, 
	which produces a near-final pair.
	(Such a near-final pair becomes available because, as sown in
	Sections~\ref{subsec: proof for n greater-equal than 7}	and~\ref{subsec: proof for n equal 5},
	$E$ satisfying Invariant~($\lambda$) guarantees the existence of a pair that is either ($\lambda$)-safe
	or near-final, and as $E$ is $(\lambda)$-mixed, the later must hold.)
	We then mix this near-final pair, producing $E$ near-final 
	that can be perfectly mixed using Theorem~$\ref{thm: E near-final}$.	
	The total number of mixing operations in this sequence is at most $(8n^3s(E))^n + 144n^3s^2(E)$.
\end{proof}
